% ---------------------------------------------------------------
\section{Limitations and Future Work}
\label{sec:limitations}
% ---------------------------------------------------------------

The correctness properties established in Section~\ref{sec:properties} are conditional guarantees: each holds under a defined set of structural and deployment assumptions.  This section makes those assumptions explicit, characterizes the operational consequences of their violation, and maps each gap to a directed research agenda.  The purpose is to accurately delimit the boundary of what the current specification guarantees and what remains open.

% ---------------------------------------------------------------
\subsection{Limitations}
\label{sec:limitations:subsection}
% ---------------------------------------------------------------

\subsubsection{Human Response Latency}
\label{sec:human-response-latency}

The Bailout Protocol's Liveness property guarantees eventual human contact within a bounded wall-clock interval $T_{\max}$, but $T_{\max}$ is a function of network topology, retry parameters, and provisioned timeout constants --- not of human cognitive response time.  For operational domains with time constants measured in milliseconds --- high-frequency trading, real-time process control, autonomous vehicle navigation --- the time required for a human principal to receive, interpret, and issue a binding directive may exceed the system's operational time horizon by one to two orders of magnitude.

The protocol's prescribed failure mode under latency stress is well-defined: a node stalls in \texttt{ESCALATING} until the anchor responds or all pathways time out.  This is the intended safety-preserving behavior.  However, the operational cost of a stall is domain-dependent in ways the protocol does not currently reason about.  In a precision agriculture context, a stall that delays a corrective action by minutes is acceptable; in a financial risk-interlock context, a stall that delays a margin-call response by even seconds may expose the system to catastrophic loss.  The protocol treats the human anchor as a timing black box: there is no mechanism to distinguish whether delay reflects anchor overload, a structural mismatch between $T_{\max}$ and the anchor's operational availability, or a pathway failure.

The deeper limitation is that human response latency is provisioned statically at node setup and is not sensitive to the anchor's current operational context.  $T_{\max}$ is a deployment parameter, not an adaptive runtime parameter.  For latency-sensitive deployments, the provisioning burden falls on the system designer to establish empirically grounded latency budgets and encode them as enforceable policy.  In deployment contexts where human response time is measured in seconds rather than minutes, static $T_{\max}$ provisioning may produce operational deadlocks that the protocol cannot distinguish from adversarial suppression.

\textit{Directed research:} Empirical characterization of human response latency distributions across deployment contexts to establish statistically grounded, domain-specific latency budgets.  Integration of these budgets into the node-provisioning interface as enforceable policy parameters with runtime monitoring.  Investigation of graded escalation responses: a protocol that requests a rapid \texttt{ACK} followed by a detailed directive would better match the operational reality of human decision-making under variable load.

\subsubsection{Adversarial Conditions}
\label{sec:adversarial-conditions}

The Safety property requires that human-contact pathways cannot be permanently closed without recorded human authorization.  This guarantee holds under the assumption that the host platform provides independent scheduling priority to the Bailout Monitor (Assumption A1.1).  In an adversarial environment --- particularly one in which the host platform itself may be compromised --- this assumption may not hold.  A sophisticated attacker controlling the execution substrate can modify thread priorities, preempt the monitor thread, introduce race conditions into pathway-closure operations, or rewrite the Bailout Monitor's code before it executes.

The protocol's guarantees are platform guarantees.  They rest on invariants that the host platform must supply: non-preemptible monitor scheduling, append-only persistence for the authorization ledger, and interruptible blocking calls for pathway management.  The Bailout Protocol is designed for production AI deployments in which individual agents may behave unexpectedly but the host platform is trusted --- the standard assumption for regulated industries.  It is not designed for deployments in which the platform itself is an active adversary: national infrastructure control systems, military command-and-control environments, financial clearinghouse infrastructure.

The distinction matters for deployment scoping.  The protocol correctly prevents a misbehaving agent from suppressing its own escalation; it does not correctly prevent a compromised host from suppressing the monitor that enforces the protocol.  This gap is structural, not fixable through protocol refinement alone.

\textit{Directed research:} Formal specification of the composition interface between the Bailout Protocol and substrate-integrity mechanisms.  Investigation of whether the protocol's guarantees can be preserved under weaker platform assumptions, including partially compromised substrates.  Exploration of hardware-attested execution environments (e.g., Intel SGX, ARM TrustZone) as a deployment target for the Bailout Monitor, eliminating the OS as a trust dependency.

\subsubsection{Scalability}
\label{sec:scalability}

The protocol is specified with a 1:1 relationship between each node and its human anchor.  This model is appropriate for single-node deployments or small-scale systems, but it creates a structural bottleneck in enterprise contexts where one anchor may supervise hundreds of concurrent nodes.  Without pooling and load-balancing mechanisms, the human anchor becomes the rate-limiting step in a system whose design intent is to route exceptions away from centralized machine decision points.

The current specification treats human anchoring as a static 1:1 relationship established at node provisioning.  This model must evolve to support at minimum 1:N relationships (one anchor supervising multiple nodes with workload pooling) and, in advanced deployments, N:M relationships (anchor pools with load balancing and dynamic routing).  The correctness properties of the protocol must be shown to extend to these pooled configurations, particularly Safety: no pooling configuration may produce a state in which all human-contact pathways for a given node are permanently disabled.

Additional scalability considerations arise from spawn tree depth.  As the depth of the accountability chain grows through recursive spawning in large multi-agent systems, the time to propagate an exception to the anchor grows proportionally.  The interaction between chain depth, propagation latency, and timeout configuration must be characterized for deployments with large spawn trees.  At sufficiently large depths, the accumulated propagation delay may exceed $T_{\max}$ even under benign conditions, producing spurious pathway-exhausted resolutions that are indistinguishable from adversarial suppression in the Forensic Ledger.

\textit{Directed research:} Formal specification of human anchor pool architectures, including load-balancing policies, prioritization schemes, and routing algorithms.  Extension of Safety and Liveness proofs to pooled anchor configurations with formally verified pathway-redundancy guarantees.  Empirical evaluation of anchor pool performance under simulated high-escalation-load conditions.

\subsubsection{Exploitation Risk}
\label{sec:exploitation-risk}

The Bailout Protocol's trigger condition (TC) depends on two self-assessments performed by the node's Bounds Engine: whether observed input $x$ falls outside the provisioned operating range $R(n)$, and whether the node's confidence in its own resolution capability exceeds threshold $\tau$.  Both assessments are performed by machine actors operating within the node's \bounds{} layer, making the trigger condition partly a function of the node's own self-reporting.

This creates an exploitation surface.  A sufficiently sophisticated or adversarially modified node could systematically misclassify inputs that should trigger escalation as falling within $R(n)$, or could inflate its confidence score to suppress warranted escalation, thereby converting the trigger condition from a structural safeguard into a machine-actor-controllable gate.  The protocol's immutability constraints ($\tau$ is provisioned at initialization and cannot be adjusted at runtime by any machine actor) limit the scope of this exploitation, but they do not eliminate it: a node that is compromised at provisioning time can select an appropriately high $\tau$ value before the immutability constraint takes effect.

The Forensic Ledger provides post-hoc detectability for systematic misclassification: an auditor can reconstruct the relationship between observed input conditions and triggered escalations for each node over rolling observation windows and flag statistically anomalous patterns.  However, this detection is forensic, not preventive.  The protocol does not currently include an active anomaly detector that intercepts misclassification before it produces an unrecorded escalation.

\textit{Directed research:} Design of triggering-behavior anomaly detectors that audit the statistical relationship between observed exception conditions and triggered escalations per node over rolling observation windows.  Investigation of provenance attestation for operating-range assessments: requiring that each Bounds Engine self-assessment be counter-signed by an independent audit process, making post-hoc misclassification detectable through forensic reconstruction.  Exploration of runtime attestation that the node's $R(n)$ definition has not been altered since provisioning.

% ---------------------------------------------------------------
\subsection{Future Work}
\label{sec:future-work}
% ---------------------------------------------------------------

The following research directions extend the Bailout Protocol's guarantees to broader deployment contexts and establish a principled research program beyond the current specification.

\textbf{Formal verification.}  Full mechanical verification of the Safety, Liveness, and Accountability invariants using a formal verification tool (Coq, Isabelle/HOL, or TLA$^+$) applied to the three-layer model and the Bailout Protocol state machine.  The current paper establishes the invariants through proof sketches grounded in the deterministic transition relation; the directed research agenda targets full formal verification as a precondition for deployment in regulated industries where formal proof of correctness is a certification requirement.

\textbf{Adaptive timeout mechanisms.}  Replacement of static timeout parameters ($T_{\bounds}$, $T_{\prop}$, $T_{\cap}$, $T_{\human}$) with adaptive mechanisms that adjust to anchor workload, pathway latency, and deployment context while preserving the Safety invariant.  The adaptive protocol would maintain the guarantee that the anchor is always contacted within $T_{\max}$ while reducing unnecessary stalls in low-consequence, high-latency scenarios.  This requires a principled framework for reclassification of exceptions by severity and consequence, ensuring that adaptive adjustments do not inadvertently suppress warranted escalation for high-consequence exceptions.

\textbf{Distributed human anchor pools.}  Formal specification and empirical evaluation of anchor pool architectures that sustain scalability under cascading failure conditions without compromising the no-machine-actor-absorbs property.  The primary design challenge is ensuring that pooled configurations preserve the Safety invariant under all failure combinations: the failure of any single anchor in the pool must not produce a state in which all pathways to all remaining anchors are simultaneously exhausted.  This requires formal pathway-redundancy analysis across the pool topology and empirical validation under simulated cascading failure conditions.

\textbf{Exploitation-resistant trigger diagnostics.}  Development of lightweight, independent audit processes that observe the relationship between input conditions and trigger evaluations without access to the node's internal state.  These independent observers would detect systematic misclassification patterns in real time, producing alerts that propagate through the Forensic Ledger as anomaly events.  The goal is to move detection closer to prevention: not by modifying the trigger condition (which would introduce new exploitation surfaces), but by making misclassification operationally expensive through rapid forensic detection.

% ---------------------------------------------------------------
\section{Conclusion}
\label{sec:conclusion}
% ---------------------------------------------------------------

The Bailout Protocol contributes a single, architecturally precise idea: in any multi-agent system built on the \bxthree{} Framework, every exception condition that exceeds machine-actor provisioned authority must --- by architectural compulsion, not by machine-actor choice --- reach a named human principal who can bind it.  This is not an error-handling strategy selected at runtime.  It is a structural constraint on the architecture of delegation itself, embedded in the deterministic state machine of every \bxthree{} node and enforced by the \fact{} Layer's pre-commit hook, which admits no machine-actor override.

The contribution is operational and architectural, not merely philosophical.  The protocol delivers four concrete guarantees that existing escalation hierarchies do not provide.  First, it is \emph{mandatory} rather than optional: the state machine's transition to \texttt{ESCALATING} is triggered by the observational boundary condition (TC) and by the \fact{} Layer's timeout enforcement, not by a machine actor's willingness to surface a failure.  Second, it is \emph{non-suppressible}: once in \texttt{ESCALATING} state, the bypass mechanism eliminates every intermediate machine actor from the communication path to $h(n)$, so no machine actor can absorb, redirect, or delay the escalation.  Third, it is \emph{observable}: every protocol event --- trigger evaluations, state transitions, pathway selections, acknowledgments, and timeout conditions --- is committed to the append-only Forensic Ledger with SHA-256 hash-chain integrity, making retrospective verification of the protocol's correct execution deterministic rather than inferential.  Fourth, it is \emph{live}: the timeout architecture ($T_{\bounds}$, $T_{\prop}$, $T_{\cap}$, $T_{\human}$) defines a finite upper bound $T_{\max}$ on wall-clock time from trigger firing to human acknowledgment, subject only to the human anchor's availability and the Pathway Health Registry's pathway quality.

The protocol does not stand alone.  It is the fifth and final pillar of the \bxthree{} Framework's upstream accountability architecture.  Loop Isolation, Recursive Spawning, Spatial Firewall, and Sandbox Gate address the exception conditions anticipated at design time --- those within the machine-only execution graph and resolvable by bounded agentic reasoning.  The Bailout Protocol addresses the complementary remainder: conditions not anticipated, that exceed the operating range of any machine actor, and that cannot be resolved within the machine-only execution graph.  Together, the five pillars constitute structurally complete coverage of the exception space: every condition that can arise in a \bxthree{} node either is resolved within the machine-only execution graph by the appropriate preceding pillar, or is compelled by the Bailout Protocol to reach a human who can bind it.  The upstream accountability guarantee --- that \bxthree{} systems fail upward into human consciousness, never downward into autonomous chaos --- is unconditional precisely because this coverage is complete.

The \bxthree{} Framework's role in this contribution is structural: it provides the three-layer accountability architecture --- \purpose{} carrying intent and final authority, \bounds{} carrying bounded execution, and \fact{} carrying deterministic enforcement --- without which the Bailout Protocol's guarantees cannot be realized.  The immutable parent pointer, the non-adjustable confidence threshold $\tau$, the bypass mechanism, and the Forensic Ledger are each products of the three-layer model.  The protocol and the framework are not separable: the framework defines the structural properties that make the protocol possible; the protocol is the runtime mechanism that makes the framework's upstream accountability guarantee operative at execution time.  One cannot be understood without the other.

The limitations enumerated in Section~\ref{sec:limitations} are real and acknowledged.  Human response latency under burst loads, adversarial pathway integrity, scalability at high delegation depth, and compulsory escalation as an exploitation surface each represent a surface area that the current specification does not fully address.  Future work --- formal verification, adaptive timeout, distributed human anchor pools, and exploitation-resistant trigger diagnostics --- is targeted precisely at these surfaces, and each extension is designed to preserve the bypass property and the mandatory escalation guarantee rather than relax them.  The protocol is a foundation, not a finished edifice.  Its structural core is sound and its extensions are well-defined.  The foundation is what matters operationally today.
